\footnotesize%
\pstart%
\noindent%
\pagelabel{impulsecc_intro}%
Die folgenden Texte N.\;\ref{60142+60297} und N.\;\ref{60299} erweisen sich inhaltlich und in ihrer Entstehung als eng miteinander verbunden und werden aus diesem Grund einheitlich ediert.
Vom Text N.\;\ref{60142+60297} sind ferner zwei Zeugen (\textit{L\textsuperscript{1}} und \textit{L\textsuperscript{2}}) überliefert, wobei \textit{L\textsuperscript{2}} offenbar eine Reinschrift des Konzepts \textit{L\textsuperscript{1}} ist (siehe die Beschreibung der Zeugen, \refpassage{LH_35_10_17_009+37_05_041_ueberlief-1}{LH_35_10_17_009+37_05_041_ueberlief-2}).
Der Textträger des Konzeptes \textit{L\textsuperscript{1}} weist freilich kein Wasserzeichen, wohl aber äußere Merkmale italienischen Papiers auf und ist somit aller Wahrscheinlichkeit nach zwischen März 1689 und März 1690 beschrieben worden, als Leibniz durch Italien\indexO{Italien} reiste.
Im Textträger der Reinschrift \textit{L\textsuperscript{2}} liegt das Fragment eines Wasserzeichens vor, das möglicherweise einer Papiermühle aus Venedigs\indexO{Venedig} Umgebung zugeschrieben werden kann.
Daher dürfte auch \textit{L\textsuperscript{2}} während der Italienreise entstanden sein.
Die Aufzeichnung N.\;\ref{60299} überliefert in ihren ersten Zeilen den nahezu gleichen Text wie N.\;\ref{60142+60297} (siehe hierzu die Beschreibung des Zeugen, %
\refpassage{LH_37_05_043_ueberliefer-1}{LH_37_05_043_ueberliefer-2}).
Hinzu kommt, dass der Textträger von N.\;\ref{60299}, dem das Wasserzeichen fehlt, die gleichen äußeren Merkmale aufweist wie der Träger der Reinschrift N.\;\ref{60142+60297}, \textit{L\textsuperscript{2}}.
Daher ist auch für N.\;\ref{60299} eine Entstehung zur Zeit von Leibnizens Italienreise anzunehmen.\indexO{Italien}
Anhand der im Folgenden dargelegten formalen und inhaltlichen Aspekte lässt sich die Datierung beider Texte N.\;\ref{60142+60297} und 
N.\;\ref{60299} jedoch weiter einschränken.
(Die Frage nach der relativen Chronologie wird erst am Schluss der editorischen Erläuterung erörtert.)
\pend%
\pstart%
Bei N.\;\ref{60142+60297} lässt die Form der Überschrift (\textit{Propositio \dots Problema} [/] \textit{Impulsus Eccentrice directi effectum explicare}, \refpassage{60142+60297_titel-1}{60142+60297_titel-2}) vermuten, dass der Text ursprünglich als Beitrag zur umfangreichen Abhandlung \textit{Dynamica de potentia et legibus naturae corporeae} gedacht war.
(Die gleiche Form \textit{Propositio \textendash\ Problema} kommt dort mehrmals vor; vgl. etwa %
\textit{CFP}~II.1, S.\;44--46, 52, 74, 80, 86, 228, 236, 268, 292 nach der Fassung \bracks{C}.)\cite{01354}
Insbesondere wäre zu erwarten, dass der \glqq Lehrsatz\grqq\ N.\;\ref{60142+60297} für den Abschnitt \textit{De concursu corporum} gedacht war, der die im Konvolut LH~XXXV~11, 18~A--C des Leibniz-Nachlasses überlieferte Fassung der \textit{Dynamica}\cite{01354} abschließt (%
\textit{CFP}~I, S.\;392--445~\bracks{A}; II.2, S.\;490--550~\bracks{C}).\cite{01354}
Im genannten Abschnitt findet sich dennoch weder ein Lehrsatz, der inhaltlich dem Text N.\;\ref{60142+60297} entspricht, noch überhaupt eine thematische Behandlung des direkten exzentrischen Stoßes zweier Körper.
Die Lehrsätze 18 und 19 erörtern zwar ausführlich die Entstehung eines \textit{gyrus}, d.\,h., wie sich aus dem Stoß mehrerer Körper eine Drehung ergibt (was auch beim direkten exzentrischen Stoß zweier Körper der Fall ist); im Mittelpunkt der Betrachtung steht dort %
jedoch wesentlich die Frage, wie kreisförmige bzw. kurvige Bewegungen aus gradlinigen gleichförmigen entstehen können, wohingegen der direkte exzentrische Stoß nur marginal besprochen wird (vgl. %
\textit{CFP}~I, S.\;419--436~\bracks{A}; II.2, S.\;520--538~\bracks{C}).\cite{01354}
Dieser bemerkenswerte Tatbestand ermöglicht \textendash\ unter Berücksichtigung weiterer biographischer Umstände \textendash\ eine genauere Datierung von N.\;\ref{60142+60297}.
\pend%
\pstart%
\edlabel{impulsecc_intro_Dynam-1}%
Leibniz begann wohl im Hochsommer 1689 den Text der \textit{Dynamica}\cite{01354} zu verfassen, während er sich in Rom\indexO{Rom} aufhielt (siehe hierüber \textsc{Robinet} 1988, S.\;81\textendash83; 259\textendash262\cite{01240}).
Ein weiter Teil der überlieferten Abhandlung war offenbar bereits niedergeschrieben worden, als in Dezember 1689 der Hofgelehrte Rudolf Christian von Bodenhausen,\indexN{\textls{Bodenhausen} (Bodenus, Bodenausen), Rudolf Christian, Freiherr v. 1698}\indexN{\textls{Bodenhausen} (Bodenus, Bodenausen), Rudolf Christian, Freiherr v. 1698} den Leibniz in Florenz\indexO{Florenz} kennengelernt hatte, sich dem \glqq Projekt\grqq\ anschloss.
Ihm war \textendash\ wie seinem (ersten) Brief an Leibniz vom 31.\;Dezember 1689 zu entnehmen ist \textendash\ der bislang entstandene Teil der \textit{Dynamica}\cite{01354} anvertraut worden mit dem Auftrag, daraus eine Reinschrift für den Druck anzufertigen (siehe \textit{LSB} III,\;4 N.\;223, S.\;436.5\textendash7).\cite{01384}
Leibniz erwartete von Bodenhausen\indexN{\textls{Bodenhausen} (Bodenus, Bodenausen), Rudolf Christian, Freiherr v. 1698} auch kritische Rückmeldungen über den Text und plante, ihm zu einem späteren Zeitpunkt auch Ergänzungen und die noch fehlenden Teile der Abhandlung zu senden (siehe seinen Brief an Bodenhausen\indexN{\textls{Bodenhausen} (Bodenus, Bodenausen), Rudolf Christian, Freiherr v. 1698} vom 13.\;Januar 1690; \textit{LSB} III,\;4 N.\;227, S.\;444.18\textendash23).\cite{01385}
Bereits am 21.\;Januar \textendash\ und dann wieder am 28.\;Januar, zu Ende Februar und am 4.\;März \textendash\ meldete Bodenhausen\indexN{\textls{Bodenhausen} (Bodenus, Bodenausen), Rudolf Christian, Freiherr v. 1698} an Leibniz, er komme mit der Arbeit an den \textit{Dynamica}\cite{01354} zügig voran und warte deshalb dringend auf das Ende des Abschnitts \textit{De concursu corporum} sowie auf die die Abhandlung abschließenden \textit{Miscellanea} (siehe \textit{LSB} III,\;4 N.\;229, \textit{K\textsuperscript{1}}, S. 448.4\textendash7, \textit{K\textsuperscript{2}}, S.\;450.6\textendash10;\cite{01386} N.\;234, S.\;458.8\textendash11;\cite{01387} N.\;240, S.\;448.16\textendash461.1;\cite{01388} N.\;241, S.\;471.2\textendash3\cite{01389}).
Am 4.\;März 1690 teilte ihm Leibniz aus Venedig mit, er werde erst auf der Rückfahrt nach Deutschland\indexO{Deutschland} in der Lage sein, den Abschnitt über den Stoß abzuschließen und die noch fehlenden Teile der \textit{Dynamica}\cite{01354} anzufertigen (\textit{LSB} III,\;4 N.\;242, S.\;473.25\textendash474.1;\cite{01390} die am 11.\;März nach Florenz gesendeten Ergänzungen enthielten tatsächlich nicht den erwarteten Schluss \textit{De concursu corporum}, vgl. \textit{LSB} III,\;4 N. 243, S.\;478.23\textendash479.5\cite{01391}).
Ähnlich drückte er sich auch wenige Tage vor der Abfahrt aus Venedig\indexO{Venedig} nach Augsburg\indexO{Augsburg} aus: In seinem Brief an Bodenhausen\indexN{\textls{Bodenhausen} (Bodenus, Bodenausen), Rudolf Christian, Freiherr v. 1698} vom 18.\;März 1690, in dem auch ein neues Inhaltsverzeichnis für die Abhandlung mit gesendet wurde, kündigt Leibniz an, er werde die \glqq sectionem \textit{De concursu corporum} forderlichst absolvieren\grqq\ und dann samt neuen bemerkenswerten Ergebnissen \textendash\ etwa über den Widerstand der Luft \textendash\ zuschicken; der Versuch, die \textit{Dynamica}\cite{01354} abzuschließen, sei doch für die Fahrt nach Augsburg\indexO{Augsburg} geplant (\textit{LSB} III,\;4 N.\;245, S.\;485.15\textendash23).\cite{01392}
Dazu aber ist Leibniz offenbar nie gekommen.
Am 29.\;Juli 1690, einen guten Monat nach seiner Ankunft in Hannover,\indexO{Hannover} erinnerte ihn Bodenhausen\indexN{\textls{Bodenhausen} (Bodenus, Bodenausen), Rudolf Christian, Freiherr v. 1698} aus Florenz\indexO{Florenz} wieder daran, dass der Abschnitt über den Stoß \textendash\ unter anderem \textendash\ noch unvollständig war (\textit{LSB} III,\;4 N.\;268, S.\;539.5\textendash6).\cite{01393}
Denselben Tenor weist auch eine Anmerkung in seinem Brief an Leibniz vom 23.\;Juni 1691 auf (\textit{LSB} III,\;5 N.\;25, S.\;127.2\textendash3).\cite{01394}
Und noch am 17.\;November 1694 ermunterte Bodenhausen\indexN{\textls{Bodenhausen} (Bodenus, Bodenausen), Rudolf Christian, Freiherr v. 1698} mit sanften Tönen Leibniz dazu, den Abschnitt \textit{De concursu corporum} abzuschließen (\textit{LSB} III,\;6 N.\;71, S~210.22\textendash24).\cite{01395}
Vergeblich.
Seit Anfang Juli 1690 hat ihm Leibniz keine neuen Textteile der \textit{Dynamica}\cite{01354} mehr zugeschickt und die Abhandlung in seinen Briefen an den Florentiner\indexO{Florenz} Hofgelehrten kaum noch erwähnt.
In seinem Schreiben an Leibniz vom 8.\;Februar 1695 nahm Bodenhausen\indexN{\textls{Bodenhausen} (Bodenus, Bodenausen), Rudolf Christian, Freiherr v. 1698} mit leicht verschleierter Enttäuschung vom gemeinsamen Projekt Abstand (\textit{LSB} III,\;6 N.\;93, S.\;279.16\textendash280.11).\cite{01396}
Damit war die Zusammenarbeit beider Gelehrter an den \textit{Dynamica}\cite{01354} beendet und das Schicksal der Abhandlung besiegelt.%
\edlabel{impulsecc_intro_Dynam-2}
\pend%
\pstart%
Der Text N.\;\ref{60142+60297} gehört offenbar nicht zu dem Teil der \textit{Dynamica},\cite{01354} den Leibniz seinem Mitstreiter Bodenhausen\indexN{\textls{Bodenhausen} (Bodenus, Bodenausen), Rudolf Christian, Freiherr v. 1698} anvertraute und dieser ordnete, prüfte und abschrieb (das ganze Konvolut, heute im Nachlass unter Signatur LH~XXXV~11, 18A--C aufbewahrt, blieb bei Bodenhausen\indexN{\textls{Bodenhausen} (Bodenus, Bodenausen), Rudolf Christian, Freiherr v. 1698} in Florenz\indexO{Florenz} bis zu dessen Tode im April 1698 und wurde anschließend von Leibniz nach Hannover\indexO{Hannover} zurückgefordert, siehe \textsc{Robinet} 1988, S.\;262\textendash267).\cite{01240}
Die Annahme liegt somit nahe, dass N.\;\ref{60142+60297} erst nach Beginn der Zusammenarbeit beider Gelehrter entstand, d.\,h. nicht vor Dezember 1689.
Insbesondere dürfte N.\;\ref{60142+60297} für jenen Schlussteil des Abschnitts \textit{De concursu corporum} gedacht worden sein, den Leibniz, allen Ermunterungen Bodenhausens\indexN{\textls{Bodenhausen} (Bodenus, Bodenausen), Rudolf Christian, Freiherr v. 1698} zum Trotz, nie anfertigen sollte.
Dass das Konzept \textit{L\textsuperscript{1}} bereits im Herbst 1689 (in Rom\indexO{Rom} oder Florenz\indexO{Florenz}) verfasst wurde, ist nicht auszuschließen:
Leibniz könnte \textit{L\textsuperscript{1}} absichtlich bei sich behalten haben, um den Text zu prüfen oder weiterzuentwickeln.
Der inhaltliche Zusammenhang mit dem für die \textit{Dynamica} vorgesehenen, nie vervollständigten Schlussteil der Stoßlehre lässt eine Entstehung vor Dezember 1689 jedoch als unwahrscheinlich erscheinen.
Jedenfalls sind beide Zeugen \textit{L\textsuperscript{1}} und \textit{L\textsuperscript{2}} aller Wahrscheinlichkeit nach noch vor Leibnizens Abreise aus Italien\indexO{Italien} Ende März 1690 entstanden, wie dies dem verwendeten Papier zu entnehmen ist.
Selbst wenn Leibniz beide Zeugen oder lediglich die Reinschrift \textit{L\textsuperscript{2}} erst auf der Rückfahrt nach Deutschland\indexO{Deutschland} abgefasst haben sollte, muss N.\;\ref{60142+60297} bereits vor dem auf Papier der Klostermühle zu Thierhaupten\indexO{Thierhaupten bei Augsburg} bei Augsburg\indexO{Augsburg} überlieferten Entwurf N.\;\ref{60300} vorgelegen haben, d.\,h. höchstwahrscheinlich noch im April 1690, allerspätestens aber vor dem folgenden Herbst (siehe die Datierungsgründe zu N.\;\ref{60300}, S.\;\pageref{LH_37_05_044_datierung}).
Denn in N.\;\ref{60142+60297} \textendash\ ebenso wie in der verwandten Aufzeichnung N.\;\ref{60299} \textendash\ wird behauptet, von einem direkten exzentrischen Stoß rühre eine Drehung her, deren Mittelpunkt mit dem Schwerpunkt des indirekt angestoßenen Körperteils zusammenfalle (\refpassage{60142+60297_centrumgyrationis-1}{60142+60297_centrumgyrationis-2} u. \refpassage{LH_37_05_043_043v_centrgyr-1}{LH_37_05_043_043v_urzelle-2}).
Diese Behauptung wird in N.\;\ref{60300} ausdrücklich zurückgewiesen und durch die Ansicht ersetzt, dass der Mittelpunkt der Drehung mit dem Schwerpunkt des gesamten angestoßenen Körpers zusammenfalle (\refpassage{LH_37_05_044_044r_cntrgyr-1}{LH_37_05_044_044r_cntrgyr-2}).
\pend%
\pstart%
Die Entstehungszeit der Aufzeichnung N.\;\ref{60299} \textendash\ und somit die relative Chronologie von N.\;\ref{60142+60297} und N.\;\ref{60299} \textendash\ lässt sich mit hinreichender Wahrscheinlichkeit durch einen Vergleich beider Texte ermitteln.
Ausschlaggebend ist der bereits eingangs hervorgehobene Umstand, dass in beiden Texten die Anfangspassage \textit{Directio conatuum ... quam ut} nahezu identisch ist.
Insbesondere stellt sich heraus, dass die ursprüngliche Fassung dieser Passage im Konzept N.\;\ref{60142+60297}, \textit{L\textsuperscript{1}} (\refpassage{60142+60297_anfang-1}{60142+60297_anfang-2} samt Variantenapparat) mit nur geringfügigen Abweichungen den gleichen Text wie der Anfang der Aufzeichnung N.\;\ref{60299} überliefert (\refpassage{LH_37_05_043r_anfang-1}{LH_37_05_043r_anfang-2}).
Dem Variantenapparat von N.\;\ref{60142+60297} ist ferner zu entnehmen, dass im Konzept \textit{L\textsuperscript{1}} die besagte Anfangspassage ursprünglich fehlte und erst nachträglich (am oberen Seitenrand) ergänzt wurde.
Bei der Abfassung der Reinschrift N.\;\ref{60142+60297}, \textit{L\textsuperscript{2}} wurde die Anfangspassage dann leicht geändert.
Denkbar ist folglich, dass das Konzept N.\;\ref{60142+60297}, \textit{L\textsuperscript{1}} (ohne Anfangspassage) bereits vorlag, als  die inhaltlich breitere Aufzeichnung N.\;\ref{60299} verfasst wurde, in welcher nicht so sehr der direkte exzentrische Stoß im Mittelpunkt steht, als vielmehr die allgemeinere Frage nach der \textit{directio conatuum eccentrica}.
Leibniz könnte anschließend bemerkt haben, dass die ersten Zeilen der Aufzeichnung N.\;\ref{60299} sich auch als Anfang des %
Lehrsatzes N.\;\ref{60142+60297} über den direkten exzentrischen Stoß eigneten, weshalb sie dem Konzept \textit{L\textsuperscript{1}} ergänzt wurden.
Daraufhin könnte er \textendash\ auf gleichem Papier wie bei der Aufzeichnung N.\;\ref{60299} \textendash\ die Reinschrift N.\;\ref{60142+60297}, \textit{L\textsuperscript{2}} angefertigt haben.
Trifft diese Darstellung zu, so ist berechtigt, von einer verschachtelten Entstehung der Texte N.\;\ref{60142+60297} und N.\;\ref{60299} zu reden.
Insbesondere ist eine gemeinsame Datierung der Reinschrift N.\;\ref{60142+60297},\;\textit{L\textsuperscript{2}} und der Aufzeichnung N.\;\ref{60299} naheliegend:
Während das Konzept N.\;\ref{60142+60297}, \textit{L\textsuperscript{1}} bereits früher \textendash\ wohl Ende 1689 oder Anfang 1690 in Oberitalien\indexO{Italien}, frühestens im Herbst 1689 in Rom\indexO{Rom} oder Florenz\indexO{Florenz} \textendash\ entstanden sein kann, dürften die Reinschrift N.\;\ref{60142+60297}, \textit{L\textsuperscript{2}} und die Aufzeichnung N.\;\ref{60299} erst in Februar oder März 1690 in Venedig\indexO{Venedig} oder aber etwas später in Deutschland\indexO{Deutschland} niedergeschrieben worden sein, jedenfalls vor dem Entwurf N.\;\ref{60300}, d.\,h. spätestens vor dem Herbst 1690.
Ein abweichender Entstehungsablauf und somit eine andere relative Chronologie der Texte N.\;\ref{60142+60297} und \ref{60299} sind jedoch nicht auszuschließen.
\pend%
\selectlanguage{latin}%
\frenchspacing%
\normalsize%
